\scanpage{80}
the whole structure being described canonically in terms of the bitorsor $\operatorname{Hom}_{\widehat\Pi_1(X)}(x,y)$ under $(\widehat\Pi_1(X,y),\widehat\Pi_1(X,x))$ and the action of $\Gamma$ on \deleted{all} this paraphernalia. One will recall above all that $f_{\unsure{y,x|x,y}}$ and $g_{y,x}$ determine each other, and that they are determined by the knowledge of \uline{a single} \deleted{$f_l$ \inserted{$:\Gamma\longrightarrow\widehat\Pi_1(X,x)$} (or} \deleted{$g_l$ (ie. $f_l(u)$} $f_l$, or $g_l$
\[
\tag{33}\left\{\begin{aligned}
f_l&:\Gamma\longrightarrow\widehat\Pi_1(X,x)\\
g_l&:\Gamma\longrightarrow\widehat\Pi_1(Y,y)
\end{aligned}\right.
\qquad l:x\longrightarrow y\quad(\text{in }\Pi_1(X))
\]
(if one has \deleted{\unsure{a}} a $1$-connected $T^\alpha$ $\subset X$, containing $x,y$) they are known when one knows only $g^{T^\alpha}_{x,y}$.

\margin{\uline{NB} $f_{x,y}$ and $f_{y,x}$, or $g_{x,y}$ and $g_{y,x}$, determine each other, for in terms of the path $l$, one has for example
\[
\tag{34}\left\{\begin{aligned}
g_{l^{-1}}(u)&=f_l(u)^{-1}\\
f_{l^{-1}}(u)&=g_l(u)^{-1}
\end{aligned}\right.
\]}

When $X$ embeds in an abelian variety, more generally if it embeds in an extension of an abelian variety by a torus (for example $X$ an ``anabelian'' algebraic curve), then the Mordell-Weil theorem implies that for a fixed $x_0$ \inserted{$\in X_K$} (as a reference point --- but this $x_0$ \deleted{\unsure{comes}} \unsure{for|plays} a bogus role!)
\margin{79}
